% Options for packages loaded elsewhere
\PassOptionsToPackage{unicode}{hyperref}
\PassOptionsToPackage{hyphens}{url}
\PassOptionsToPackage{dvipsnames,svgnames,x11names}{xcolor}
\documentclass[
  11pt,
]{article}
\usepackage{xcolor}
\usepackage[margin=1in]{geometry}
\usepackage{amsmath,amssymb}
\setcounter{secnumdepth}{-\maxdimen} % remove section numbering
\usepackage{iftex}
\ifPDFTeX
  \usepackage[T1]{fontenc}
  \usepackage[utf8]{inputenc}
  \usepackage{textcomp} % provide euro and other symbols
\else % if luatex or xetex
  \usepackage{unicode-math} % this also loads fontspec
  \defaultfontfeatures{Scale=MatchLowercase}
  \defaultfontfeatures[\rmfamily]{Ligatures=TeX,Scale=1}
\fi
\usepackage{lmodern}
\ifPDFTeX\else
  % xetex/luatex font selection
\fi
% Use upquote if available, for straight quotes in verbatim environments
\IfFileExists{upquote.sty}{\usepackage{upquote}}{}
\IfFileExists{microtype.sty}{% use microtype if available
  \usepackage[]{microtype}
  \UseMicrotypeSet[protrusion]{basicmath} % disable protrusion for tt fonts
}{}
\makeatletter
\@ifundefined{KOMAClassName}{% if non-KOMA class
  \IfFileExists{parskip.sty}{%
    \usepackage{parskip}
  }{% else
    \setlength{\parindent}{0pt}
    \setlength{\parskip}{6pt plus 2pt minus 1pt}}
}{% if KOMA class
  \KOMAoptions{parskip=half}}
\makeatother
\usepackage{listings}
\newcommand{\passthrough}[1]{#1}
\lstset{defaultdialect=[5.3]Lua}
\lstset{defaultdialect=[x86masm]Assembler}
\usepackage{longtable,booktabs,array}
\usepackage{caption}
\captionsetup[table]{skip=6pt}
\newcounter{none} % for unnumbered tables
\usepackage{calc} % for calculating minipage widths
% Correct order of tables after \paragraph or \subparagraph
\usepackage{etoolbox}
\makeatletter
\patchcmd\longtable{\par}{\if@noskipsec\mbox{}\fi\par}{}{}
\makeatother
% Allow footnotes in longtable head/foot
\IfFileExists{footnotehyper.sty}{\usepackage{footnotehyper}}{\usepackage{footnote}}
\makesavenoteenv{longtable}
\setlength{\emergencystretch}{3em} % prevent overfull lines
\providecommand{\tightlist}{%
  \setlength{\itemsep}{0pt}\setlength{\parskip}{0pt}}
\usepackage{bookmark}
\IfFileExists{xurl.sty}{\usepackage{xurl}}{} % add URL line breaks if available
\urlstyle{same}
% fallback for those not using the hyperref driver hyperxmp:
\makeatletter
\@ifundefined{xmpquote}{\newcommand{\xmpquote}[1]{#1}}{}
\makeatother
\hypersetup{
  pdftitle={Declarative Handler Topology Synthesis for Arrow Effects},
  pdfauthor={Luis Emidio},
  colorlinks=true,
  linkcolor={blue},
  filecolor={Maroon},
  citecolor={blue},
  urlcolor={blue},
  pdfcreator={LaTeX via pandoc}}

\title{Declarative Handler Topology Synthesis for Arrow Effects}
\author{Luis Emidio}
\date{2025-08-23}

\begin{document}
\maketitle

\section{Abstract}\label{abstract}

Algebraic effects and handlers provide a modular mechanism for
separating the declaration of computational operations from their
interpretation. In Sanada's calculus of algebraic effects and handlers
for arrows, handler composition is represented explicitly by nested
\passthrough{\lstinline!handle!} expressions. This representation
exposes an important practical question: can handler topology be
specified declaratively and synthesized statically, rather than written
manually as nested syntax?

We present a declarative handler-topology system layered above Sanada's
arrow calculus. A declarative specification describes a finite set of
handlers together with precedence constraints. A compile-time synthesis
pass converts this specification into a linearly nested sequence of
\passthrough{\lstinline!handle!} expressions accepted by the underlying
arrow calculus. The synthesis problem is therefore formulated as finding
a linear extension of a handler-precedence partial order, rather than
constructing a branching runtime computation.

This distinction is important because the expressiveness limitation
identified by Sanada concerns effectful command branching within the
arrow semantics. Our synthesis mechanism introduces no such runtime
branching: it generates a fixed handler topology before execution. We
formalize the synthesis transformation, characterize its well-formedness
and precedence-preservation properties, and describe an implementation
strategy for arrow-effect languages that lowers declarative
specifications directly to existing handler-expression AST nodes. The
resulting system requires no new runtime mechanism and preserves the
semantic boundary imposed by the underlying arrow calculus.

\section{1. Introduction}\label{introduction}

Algebraic effects allow programs to describe operations independently of
their interpretation. Effect handlers subsequently determine how those
operations are handled. This separation provides a useful foundation for
modular programming, but the composition of handlers remains an explicit
part of the program structure.

In Sanada's calculus for algebraic effects and handlers for arrows,
handler composition is expressed by nested
\passthrough{\lstinline!handle!} constructs:

\begin{lstlisting}
handle (handle P with H₁) with H₂
\end{lstlisting}

The nesting determines handler scope and therefore potentially affects
program behavior. The programmer must consequently encode the desired
handler topology directly in the term.

We investigate whether this topology can instead be specified
declaratively.

The problem considered in this paper is narrower than general topology
synthesis. We do not attempt to infer arbitrary computational graphs,
physical layouts, or dynamic effect-routing structures. Rather, we
consider the specific problem of \textbf{synthesizing the static nesting
topology of arrow effect handlers}.

Our central observation is that the synthesis problem is naturally
expressed as an ordering problem. Given a finite collection of handlers
and constraints describing their relative precedence, the compiler can
compute a linear extension of those constraints and generate the
corresponding nested \passthrough{\lstinline!handle!} expression.

This approach is compatible with the structural restriction of the arrow
calculus. In the semantic setting considered by Sanada, the algebraic
structure associated with arrows corresponds to a sequence of operations
rather than the branching operation trees associated with ordinary
monads. In particular, the problematic case is runtime branching between
effectful commands. Our synthesis pass does not introduce such
branching: its output is a statically determined sequence of nested
handlers.

The contributions of this paper are:

\begin{enumerate}
\def\labelenumi{\arabic{enumi}.}
\tightlist
\item
  \textbf{A declarative handler-topology model.} We define handler
  declarations independently of their syntactic nesting and express
  precedence as a relation over handlers.
\item
  \textbf{A static synthesis algorithm.} We reduce handler-topology
  synthesis to finding a linear extension of the declared precedence
  relation and lower the result to nested
  \passthrough{\lstinline!handle!} expressions.
\item
  \textbf{Structural correctness properties.} We establish that
  synthesis produces valid arrow-calculus terms and preserves declared
  handler precedence.
\item
  \textbf{A minimal implementation.} We describe a compile-time AST
  transformation for arrow-effect languages that reuses existing handler
  representations and virtual machines.
\item
  \textbf{An explicit expressiveness boundary.} We distinguish static
  synthesis, which produces a fixed handler topology, from runtime
  dynamic selection of effectful computations, which lies outside the
  arrow semantics considered here.
\item
  \textbf{A Lean 4 formalization.} The correctness properties are
  mechanically verified in Lean 4 with no admitted proofs (§7).
\end{enumerate}

\section{2. Background}\label{background}

\subsection{2.1 Sanada's Arrow Calculus}\label{sanadas-arrow-calculus}

Sanada extends the arrow calculus of Lindley, Wadler, and Yallop with
operations and handlers.

The relevant syntax (Figure 2 of the thesis) is:

\begin{lstlisting}
Terms       M, N, L ::= x | ⟨M, N⟩ | fst M | snd M | λx:A. M | M N | λ•x:A. P
Values      V, W    ::= x | ⟨V, W⟩ | λx:A. M | λ•x:A. P
Commands    P, Q, R ::= ⌊M⌋ | let x ⇐ P in Q | L • M | op(M) | handle R with H
Handlers    H       ::= { x:C → P } ∪ { op, k:δ_D z:γ → Q_op }_{op ∈ Σ}
\end{lstlisting}

A handler \passthrough{\lstinline!H!} consists of a return clause
\passthrough{\lstinline!\{ x:C → P \}!} and a set of operation clauses
\passthrough{\lstinline!\{ op, k:δ\_D z:γ → Q\_op \}!}. Each operation
clause binds:

\begin{itemize}
\tightlist
\item
  \passthrough{\lstinline!z : γ!} --- the operation argument (arity γ);
\item
  \passthrough{\lstinline!k : δ\_D!} --- the continuation (coarity δ);
\item
  \passthrough{\lstinline!Q\_op!} --- the handler body for that
  operation.
\end{itemize}

The construct

\begin{lstlisting}
handle R with H
\end{lstlisting}

evaluates command \passthrough{\lstinline!R!} under handler
\passthrough{\lstinline!H!}.

Handler nesting determines which handler is considered first when an
operation is encountered. Thus the following two terms are structurally
different:

\begin{lstlisting}
handle (handle P with H₁) with H₂
\end{lstlisting}

and:

\begin{lstlisting}
handle (handle P with H₂) with H₁
\end{lstlisting}

A declarative system must therefore synthesize this ordering explicitly.

\subsection{2.2 Handler Topology}\label{handler-topology}

We distinguish handler topology from computational topology.

\textbf{Handler topology} is the nesting relation between effect
handlers:

\begin{lstlisting}
H₁
 └── H₂
      └── H₃
           └── P
\end{lstlisting}

\textbf{Computational topology}, by contrast, describes dependencies
between computations.

The present paper concerns only the former.

This distinction is important because the system does not infer a
dataflow graph or physical layout. Its output is a valid handler nesting
that realizes a declaratively specified precedence relation.

\subsection{2.3 Arrow Operations and Sequential
Structure}\label{arrow-operations-and-sequential-structure}

Sanada's semantic treatment of arrow effects is based on strong
promonads. A relevant consequence is that the algebraic structure
associated with arrows represents a sequence of operations rather than
the branching operation trees associated with ordinary monadic effects.

This restriction has two consequences for our system.

First, a declarative handler topology must ultimately produce a linear
nesting of handlers.

Second, a runtime construct that dynamically chooses between effectful
commands based on the result of another effect cannot simply be
introduced into the arrow calculus without changing its semantic model.

Our synthesis mechanism avoids both issues by computing the handler
ordering statically.

\section{3. The Declarative Handler
Model}\label{the-declarative-handler-model}

\subsection{3.1 Handler Declarations}\label{handler-declarations}

Let a declarative handler specification be:

\begin{lstlisting}
D = (H, ≺)
\end{lstlisting}

where:

\begin{itemize}
\tightlist
\item
  \passthrough{\lstinline!H!} is a finite set of handler declarations;
\item
  \passthrough{\lstinline!≺!} is a precedence relation over
  \passthrough{\lstinline!H!}.
\end{itemize}

For example:

\begin{lstlisting}
fetch ≺ cache
cache ≺ log
\end{lstlisting}

states that \passthrough{\lstinline!fetch!} must be nested inside
\passthrough{\lstinline!cache!}, and \passthrough{\lstinline!cache!}
must be nested inside \passthrough{\lstinline!log!}.

The relation is not itself a nesting tree. It describes constraints on
the eventual linear handler topology.

\subsection{3.2 Partial Orders}\label{partial-orders}

The precedence relation is assumed to form a partial order.

Thus it must be:

\begin{itemize}
\tightlist
\item
  irreflexive;
\item
  transitive;
\item
  asymmetric.
\end{itemize}

If two handlers are unrelated by the declared relation, the compiler is
free to choose either relative order, subject to any additional semantic
requirements imposed by the implementation.

If the relation contains a cycle:

\begin{lstlisting}
H₁ ≺ H₂
H₂ ≺ H₃
H₃ ≺ H₁
\end{lstlisting}

then no valid linear handler topology exists.

The compiler must reject such a specification rather than arbitrarily
choosing an order.

\subsection{3.3 Linear Extensions}\label{linear-extensions}

A \textbf{linear extension} of \passthrough{\lstinline!(H, ≺)!} is a
sequence:

\begin{lstlisting}
H₁, H₂, ..., Hₙ
\end{lstlisting}

such that:

\begin{lstlisting}
Hᵢ ≺ Hⱼ ⇒ i < j
\end{lstlisting}

The synthesis problem is therefore:

\begin{quote}
Given a finite declarative handler specification
\passthrough{\lstinline!(H, ≺)!}, compute a linear extension of its
precedence relation.
\end{quote}

A standard topological-sort algorithm is sufficient for this purpose.

The important point is that this is a compile-time ordering problem. No
runtime branching is necessary.

\section{4. Handler Topology
Synthesis}\label{handler-topology-synthesis}

\subsection{4.1 Synthesis Function}\label{synthesis-function}

Let:

\begin{lstlisting}
S(D, P)
\end{lstlisting}

be the synthesis function that takes a declarative specification
\passthrough{\lstinline!D!} and an arrow command
\passthrough{\lstinline!P!}.

Suppose the linear extension computed from \passthrough{\lstinline!D!}
is:

\begin{lstlisting}
H₁, H₂, ..., Hₙ
\end{lstlisting}

The synthesized command is:

\begin{lstlisting}
handle (... handle (handle P with H₁) with H₂ ... ) with Hₙ
\end{lstlisting}

Equivalently:

\begin{lstlisting}
S(D, P) =
  handle (...(handle (handle P with H₁) with H₂)... ) with Hₙ
\end{lstlisting}

Thus the first handler in the linear extension is the innermost handler.

\subsection{4.2 Example}\label{example}

Consider:

\begin{lstlisting}
events:

  - handler: FetchHandler
    handles: [fetch]
    priority: 10

  - handler: CacheHandler
    handles: [cache]
    priority: 20

  - handler: LogHandler
    handles: [log]
    priority: 30

body: ...
\end{lstlisting}

The priority declarations induce:

\begin{lstlisting}
FetchHandler ≺ CacheHandler ≺ LogHandler
\end{lstlisting}

The synthesizer produces:

\begin{lstlisting}
handle (
  handle (
    handle body with FetchHandler
  ) with CacheHandler
) with LogHandler
\end{lstlisting}

The declarative representation therefore describes the ordering
constraint without requiring the programmer to manually encode the
nesting.

\subsection{4.3 Priority as a Convenience
Syntax}\label{priority-as-a-convenience-syntax}

A total numeric priority is one convenient way to specify a precedence
relation:

\begin{lstlisting}
priority(H₁) < priority(H₂)
\end{lstlisting}

implies:

\begin{lstlisting}
H₁ ≺ H₂
\end{lstlisting}

However, priority values are not fundamental to the model.

The more general representation is a partial-order specification:

\begin{lstlisting}
H₁ ≺ H₂
H₁ ≺ H₃
\end{lstlisting}

which leaves \passthrough{\lstinline!H₂!} and
\passthrough{\lstinline!H₃!} unordered.

This distinction allows the compiler to reject genuinely contradictory
constraints while avoiding unnecessary ordering requirements.

\section{5. Synthesis Algorithm}\label{synthesis-algorithm}

\subsection{5.1 Algorithm}\label{algorithm}

Given:

\begin{lstlisting}
D = (H, ≺)
\end{lstlisting}

and body \passthrough{\lstinline!P!}:

\begin{enumerate}
\def\labelenumi{\arabic{enumi}.}
\item
  Validate that \passthrough{\lstinline!≺!} is acyclic.
\item
  Compute a linear extension of \passthrough{\lstinline!≺!}.
\item
  Initialize \passthrough{\lstinline!R = P!}.
\item
  For each handler in the linear extension, from innermost to outermost:

  \begin{itemize}
  \tightlist
  \item
    construct a \passthrough{\lstinline!handle!} node containing
    \passthrough{\lstinline!R!};
  \item
    replace \passthrough{\lstinline!R!} with the new node.
  \end{itemize}
\item
  Return \passthrough{\lstinline!R!}.
\end{enumerate}

In pseudocode:

\begin{lstlisting}
synthesize(D, P):
    ordered = topological_sort(D.precedence)

    result = P

    for H in ordered:
        result = handle(result, H)

    return result
\end{lstlisting}

\subsection{5.2 Complexity}\label{complexity}

Let:

\begin{itemize}
\tightlist
\item
  \passthrough{\lstinline!n = |H|!} be the number of handlers;
\item
  \passthrough{\lstinline!m = |≺|!} be the number of precedence
  constraints.
\end{itemize}

Using a standard topological-sort algorithm, ordering requires:

\begin{lstlisting}
O(n + m)
\end{lstlisting}

time.

Constructing the resulting nested AST requires:

\begin{lstlisting}
O(n)
\end{lstlisting}

additional time.

Thus the synthesis phase is linear in the size of the handler
specification, assuming the precedence relation is represented
explicitly.

\subsection{5.3 Determinism and Multiple Linear
Extensions}\label{determinism-and-multiple-linear-extensions}

A partial order may have multiple valid linear extensions. For example,
given:

\begin{lstlisting}
A ≺ C
B ≺ C
\end{lstlisting}

both \passthrough{\lstinline![A, B, C]!} and
\passthrough{\lstinline![B, A, C]!} are valid linear extensions. The
choice between them affects which of \passthrough{\lstinline!A!} or
\passthrough{\lstinline!B!} is the innermost handler, which may in turn
affect runtime behavior if \passthrough{\lstinline!A!} and
\passthrough{\lstinline!B!} intercept overlapping operations.

For reproducible builds, the compiler can impose a deterministic
tie-breaking rule, such as declaration order or handler identifier.

This does not change the declared precedence relation:

\begin{lstlisting}
H₁ ≺ H₂
\end{lstlisting}

remains mandatory, while unrelated handlers receive a deterministic but
otherwise arbitrary ordering.

The question of when two handlers can be safely reordered without
changing behavior is a semantic question (see §14.1).

\section{6. Correctness}\label{correctness}

The central correctness properties concern the relationship between the
declarative specification and the generated syntax.

\subsection{6.1 Well-Formedness}\label{well-formedness}

We define a command as \textbf{well-formed} if every handler in it has
at least one operation clause and all sub-terms are well-formed. This is
a structural property checked recursively:

\begin{lstlisting}
isWellFormed(handle P with H) = |H.clauses| > 0 ∧ isWellFormed(P)
isWellFormed(let x ⇐ P in Q)   = isWellFormed(P) ∧ isWellFormed(Q)
isWellFormed(cond P Q R)       = isWellFormed(P) ∧ isWellFormed(Q) ∧ isWellFormed(R)
isWellFormed(op(M))            = true
isWellFormed(⌊M⌋)              = true
\end{lstlisting}

\textbf{Proposition 1 --- Well-Formedness.}

For every finite acyclic handler specification
\passthrough{\lstinline!D = (H, ≺)!} and well-formed command
\passthrough{\lstinline!P!}, if every handler in the linear extension of
\passthrough{\lstinline!≺!} has at least one clause, then
\passthrough{\lstinline!S(D, P)!} is well-formed.

\textbf{Proof sketch.}

The synthesis folds \passthrough{\lstinline!handle!} over the linear
extension. Each step wraps the current term with a handler
\passthrough{\lstinline!Hᵢ!} that has at least one clause. By the
definition of \passthrough{\lstinline!isWellFormed!}, wrapping a
well-formed term with a handler that has clauses preserves
well-formedness. Induction over the handler sequence establishes the
result. The accompanying Lean formalization proves this as
\passthrough{\lstinline!synthesize\_well\_formed!}.

\subsection{6.2 Precedence Preservation}\label{precedence-preservation}

\textbf{Proposition 2 --- Precedence Preservation.}

For every pair of handlers \passthrough{\lstinline!Hᵢ!} and
\passthrough{\lstinline!Hⱼ!}:

\begin{lstlisting}
Hᵢ ≺ Hⱼ
\end{lstlisting}

implies that \passthrough{\lstinline!Hᵢ!} is nested inside
\passthrough{\lstinline!Hⱼ!} in \passthrough{\lstinline!S(D, P)!}.

\textbf{Proof sketch.}

The synthesis algorithm uses a linear extension of
\passthrough{\lstinline!≺!}. By definition of a linear extension:

\begin{lstlisting}
Hᵢ ≺ Hⱼ ⇒ position(Hᵢ) < position(Hⱼ)
\end{lstlisting}

Since handlers are folded from the innermost position toward the
outermost position, \passthrough{\lstinline!Hᵢ!} is constructed before
\passthrough{\lstinline!Hⱼ!} and is therefore nested inside it.

The key invariant is that if \passthrough{\lstinline!Hᵢ!} appears before
\passthrough{\lstinline!Hⱼ!} in the ordered list, then for any body
term, \passthrough{\lstinline!Hᵢ!} is nested inside
\passthrough{\lstinline!Hⱼ!} in the result of folding
\passthrough{\lstinline!handle!} over that list. This is proved by
induction on the \passthrough{\lstinline!AppearsBefore!} relation, with
two helper lemmas handling the fold accumulation pattern. The
accompanying Lean formalization proves this as
\passthrough{\lstinline!synthesize\_preserves\_precedence!}.

\textbf{Corollary --- Completeness.}

If \passthrough{\lstinline!ordered!} is a linear extension of
\passthrough{\lstinline!≺!} (i.e.,
\passthrough{\lstinline!∀ a b, a ≺ b ⇒ AppearsBefore(a, b, ordered)!}),
then for every constraint \passthrough{\lstinline!Hᵢ ≺ Hⱼ!} in
\passthrough{\lstinline!≺!}, \passthrough{\lstinline!Hᵢ!} is nested
inside \passthrough{\lstinline!Hⱼ!} in
\passthrough{\lstinline!synthesize(P, ordered)!}.

This follows directly from Proposition 2: the linear extension property
ensures \passthrough{\lstinline!AppearsBefore(Hᵢ, Hⱼ, ordered)!}, and
Proposition 2 converts this into structural nesting. The Lean
formalization proves this as
\passthrough{\lstinline!synthesize\_completeness!}.

The soundness of the topological sort --- the fact that it produces a
list satisfying the linear extension property --- is also formally
verified, as \passthrough{\lstinline!exists\_linear\_extension!} in
\passthrough{\lstinline!TopoSort.lean!} (§7.2). This theorem proves that
for any finite list of handlers with an irreflexive and transitive
precedence relation, there exists a linear extension realizing all
precedence constraints as \passthrough{\lstinline!AppearsBefore!}.

\subsection{6.3 No Introduced Runtime
Branching}\label{no-introduced-runtime-branching}

We define a command as \textbf{static} if it contains no
\passthrough{\lstinline!cond!} (runtime conditional branching)
constructor. The \passthrough{\lstinline!isStatic!} predicate is checked
recursively:

\begin{lstlisting}
isStatic(handle P with H) = isStatic(P)
isStatic(cond P Q R)      = false
isStatic(let x ⇐ P in Q)  = isStatic(P) ∧ isStatic(Q)
isStatic(op(M))           = true
isStatic(⌊M⌋)             = true
\end{lstlisting}

\textbf{Proposition 3 --- Static Topology.}

If the body command \passthrough{\lstinline!P!} is static (contains no
\passthrough{\lstinline!cond!}), then \passthrough{\lstinline!S(D, P)!}
is also static. Synthesis introduces no runtime branching over effectful
commands.

\textbf{Proof sketch.}

The synthesis function folds \passthrough{\lstinline!handle!} over the
ordered list. Since
\passthrough{\lstinline!isStatic(handle P with H) = isStatic(P)!},
wrapping with a handler preserves staticity. By induction over the
handler list, if the initial body is static, the result remains static.
The \passthrough{\lstinline!cond!} constructor is never introduced by
synthesis --- it only appears if it was already present in the body.

The accompanying Lean formalization proves this as
\passthrough{\lstinline!synthesize\_static\_topology!}, with the
\passthrough{\lstinline!cond!} constructor explicitly included in the
\passthrough{\lstinline!Command!} type to make the theorem non-vacuous.

\subsection{6.4 Semantic Boundary}\label{semantic-boundary}

The preceding property should not be interpreted as proving that
arbitrary dynamic effect-routing programs can be represented in the
arrow calculus.

Rather, the result establishes a narrower fact:

\begin{quote}
A static declarative description of handler precedence can be compiled
into the existing static handler-nesting mechanism without introducing
the dynamic command branching excluded by the underlying arrow model.
\end{quote}

This distinction is central to the scope of the system.

\section{7. Formalization in Lean 4}\label{formalization-in-lean-4}

The correctness properties of §6 have been mechanically verified in Lean
4. The formalization is split across the following files under
\passthrough{\lstinline!Dts/!}:

\begin{itemize}
\tightlist
\item
  \passthrough{\lstinline!Basic.lean!} --- core types
  (\passthrough{\lstinline!Command!}, \passthrough{\lstinline!Handler!},
  \passthrough{\lstinline!isStatic!},
  \passthrough{\lstinline!isWellFormed!})
\item
  \passthrough{\lstinline!Precedence.lean!} --- precedence model,
  partial orders, acyclicity
\item
  \passthrough{\lstinline!Synthesis.lean!} --- synthesis function,
  well-formedness, precedence preservation, staticity, completeness
\item
  \passthrough{\lstinline!TopoSort.lean!} --- topological sort soundness
  (\passthrough{\lstinline!exists\_linear\_extension!})
\item
  \passthrough{\lstinline!Semantics.lean!} --- trace semantics for
  \passthrough{\lstinline!Command!} (observable handler-dispatch events)
\item
  \passthrough{\lstinline!SemanticEquivalence.lean!} --- semantic
  equivalence of synthesized and manually nested terms
\item
  \passthrough{\lstinline!Commutativity.lean!} --- semantic
  commutativity of disjoint handlers, with counterexample
\end{itemize}

The root \passthrough{\lstinline!Dts.lean!} imports all modules. The
project compiles with no \passthrough{\lstinline!sorry!} (no admitted
proofs). Mathlib is used for tactic utilities only (e.g.,
\passthrough{\lstinline!omega!}, \passthrough{\lstinline!push\_neg!});
no mathematical content is imported from it.

\subsection{7.1 Model}\label{model}

The Lean model mirrors the paper's definitions:

\begin{itemize}
\tightlist
\item
  \passthrough{\lstinline!Command!} is an inductive type with
  constructors \passthrough{\lstinline!pure!},
  \passthrough{\lstinline!op!}, \passthrough{\lstinline!letin!},
  \passthrough{\lstinline!handle!}, and \passthrough{\lstinline!cond!}.
  The \passthrough{\lstinline!cond!} constructor represents runtime
  conditional branching; its inclusion makes the staticity theorem
  (Proposition 3) non-vacuous.
\item
  \passthrough{\lstinline!Precedence!} is a function
  \passthrough{\lstinline!Handler → Handler → Bool!}.
\item
  \passthrough{\lstinline!synthesize body ordered!} is a left fold over
  the ordered handler list, matching the paper's
  \passthrough{\lstinline!S(D, P)!}.
\item
  \passthrough{\lstinline!AppearsBefore!} is an inductive predicate
  capturing list ordering without relying on numeric indices.
\item
  \passthrough{\lstinline!NestedInside!} is an inductive predicate
  capturing handler nesting structure.
\end{itemize}

\subsection{7.2 Proven Theorems}\label{proven-theorems}

{\def\LTcaptype{none} % do not increment counter
\begin{longtable}[]{@{}
  >{\raggedright\arraybackslash}p{(\linewidth - 2\tabcolsep) * \real{0.5000}}
  >{\raggedright\arraybackslash}p{(\linewidth - 2\tabcolsep) * \real{0.5000}}@{}}
\toprule\noalign{}
\begin{minipage}[b]{\linewidth}\raggedright
Lean theorem
\end{minipage} & \begin{minipage}[b]{\linewidth}\raggedright
Paper proposition
\end{minipage} \\
\midrule\noalign{}
\endhead
\bottomrule\noalign{}
\endlastfoot
\passthrough{\lstinline!irreflexive\_transitive\_implies\_asymmetric!} &
Lemma: asymmetry is derivable (not assumed) \\
\passthrough{\lstinline!partialOrder\_acyclic!} & Acyclicity of partial
orders \\
\passthrough{\lstinline!synthesize\_well\_formed!} & Proposition 1
(conditional on non-empty clauses) \\
\passthrough{\lstinline!synthesize\_preserves\_precedence!} &
Proposition 2 \\
\passthrough{\lstinline!synthesize\_static\_topology!} & Proposition 3
(conditional on body being static) \\
\passthrough{\lstinline!synthesize\_completeness!} & Corollary
(completeness given linear extension) \\
\passthrough{\lstinline!reachable\_implies\_edge!} & Lemma: reachability
implies a direct edge (used by acyclicity) \\
\passthrough{\lstinline!exists\_linear\_extension!} & Topological sort
soundness (existence of linear extension) \\
\passthrough{\lstinline!synthesize\_eq\_manuallyNested!} & Syntactic
equality of synthesis and explicit nesting \\
\passthrough{\lstinline!synthesize\_equiv\_manuallyNested!} & Semantic
equivalence (Corollary: equal syntax ⇒ equal traces) \\
\passthrough{\lstinline!synthesize\_commutative\_two!} & Semantic
commutativity for disjoint handlers (§14.1) \\
\end{longtable}
}

\subsection{7.3 Assumptions}\label{assumptions}

The formalization makes no unproven assumptions. The linear extension
property used by \passthrough{\lstinline!synthesize\_completeness!} is
established by \passthrough{\lstinline!exists\_linear\_extension!} in
\passthrough{\lstinline!TopoSort.lean!}, which proves that for any
finite list of handlers with an irreflexive and transitive precedence
relation, a valid linear extension exists. This is a constructive
version of a standard result in order theory (Szpilrajn's extension
theorem).

\section{8. Relationship to Arrow
Expressiveness}\label{relationship-to-arrow-expressiveness}

\subsection{8.1 Dynamic Command
Branching}\label{dynamic-command-branching}

Consider a hypothetical program:

\begin{lstlisting}
v ⇐ op()
if v then P₁ else P₂
\end{lstlisting}

where \passthrough{\lstinline!P₁!} and \passthrough{\lstinline!P₂!} are
effectful commands.

The branch selects a future effectful computation based on the result of
an operation. This is fundamentally different from static handler
synthesis.

Our system instead performs:

\begin{lstlisting}
compile time:
    determine H₁ ≺ H₂ ≺ H₃

execution time:
    execute fixed nesting
\end{lstlisting}

The handler topology is therefore independent of runtime operation
results.

\subsection{8.2 Pure Term-Level
Branching}\label{pure-term-level-branching}

The restriction does not imply that all conditional behavior disappears.

A pure value can still participate in ordinary term-level computation.
The relevant distinction is whether the branch dynamically selects
between effectful command structures.

Our compiler does not need to alter this distinction.

\subsection{8.3 Why the Restriction Is
Useful}\label{why-the-restriction-is-useful}

The same restriction that limits dynamic expressiveness also simplifies
static synthesis.

If arbitrary runtime command branching were permitted, the handler
topology might itself depend on runtime execution paths. A compiler
would then potentially need to synthesize a branching handler structure.

Under the arrow model considered here, static handler topology remains a
linear structure:

\begin{lstlisting}
H₁
 ↓
H₂
 ↓
H₃
 ↓
P
\end{lstlisting}

This makes compile-time synthesis particularly simple.

\section{9. Declarative Language
Design}\label{declarative-language-design}

\subsection{9.1 Syntax}\label{syntax}

A concrete configuration syntax can expose handler precedence without
exposing the underlying nesting.

For example:

\begin{lstlisting}
events:
  - handler: Cache
    before: [Logging]

  - handler: Authentication
    before: [Cache]

  - handler: Logging

body:
  application
\end{lstlisting}

This corresponds to:

\begin{lstlisting}
Authentication ≺ Cache ≺ Logging
\end{lstlisting}

The configuration language need not expose
\passthrough{\lstinline!handle!} directly.

\subsection{9.2 Explicit Relations versus
Priorities}\label{explicit-relations-versus-priorities}

We recommend treating explicit precedence constraints as the semantic
core:

\begin{lstlisting}
Authentication before Cache
Cache before Logging
\end{lstlisting}

Numeric priorities can be provided as syntactic sugar:

\begin{lstlisting}
Authentication: 10
Cache: 20
Logging: 30
\end{lstlisting}

This avoids coupling the language semantics to an arbitrary numeric
scale.

\subsection{9.3 Invalid Specifications}\label{invalid-specifications}

The compiler should diagnose at least:

\begin{itemize}
\tightlist
\item
  cyclic precedence constraints;
\item
  duplicate handler identities;
\item
  references to undefined handlers;
\item
  malformed handler clauses;
\item
  conflicting declarations.
\end{itemize}

For example:

\begin{lstlisting}
A before B
B before A
\end{lstlisting}

must produce a compile-time error because no valid linear nesting
satisfies both constraints.

\section{10. Implementation}\label{implementation}

The proposed implementation is deliberately small.

The declarative frontend produces handler declarations and precedence
constraints. A synthesis pass then transforms the existing AST before VM
evaluation.

Conceptually:

\begin{lstlisting}
Declarative specification
          │
          ▼
 Handler declarations
          │
          ▼
 Precedence graph
          │
          ▼
 Topological sort
          │
          ▼
 Nested handler expressions
          │
          ▼
 Existing compiler / VM
\end{lstlisting}

The synthesis pass is therefore a pure AST transformation.

A simplified implementation in pseudocode:

\begin{lstlisting}
synthesize_topology(body, handlers):
    ordered = topological_sort(handlers.precedence)

    result = body

    for H in ordered:
        result = handle(result, H)

    return result
\end{lstlisting}

The implementation does not require a new runtime representation of
handlers. The generated nodes use the existing handler-expression
representation of the target language.

Likewise, the VM continues to perform ordinary handler lookup and
evaluation.

The principal implementation advantage is therefore not that handlers
become intrinsically cheaper, but that the declarative layer introduces
\textbf{no additional runtime dispatch mechanism}.

\section{11. Evaluation Plan}\label{evaluation-plan}

\subsection{11.1 Correctness of Ordering}\label{correctness-of-ordering}

Generate handler specifications with known precedence relations and
verify that the resulting AST has the corresponding nesting.

For example:

\begin{lstlisting}
A ≺ B ≺ C
\end{lstlisting}

must produce:

\begin{lstlisting}
handle(
  handle(
    handle(P, A),
    B
  ),
  C
)
\end{lstlisting}

\subsection{11.2 Detection of Invalid
Topologies}\label{detection-of-invalid-topologies}

Generate cyclic specifications and verify that synthesis rejects them at
compile time.

Examples include:

\begin{lstlisting}
A ≺ B
B ≺ A
\end{lstlisting}

and:

\begin{lstlisting}
A ≺ B
B ≺ C
C ≺ A
\end{lstlisting}

\subsection{11.3 Behavioral Equivalence}\label{behavioral-equivalence}

For specifications whose nesting is known manually, compare:

\begin{lstlisting}
manually nested program
\end{lstlisting}

with:

\begin{lstlisting}
declaratively synthesized program
\end{lstlisting}

and verify equivalent behavior under the existing VM.

This experiment is more significant than simply checking the AST shape
because handler nesting can affect which operation clause receives an
operation.

\subsection{11.4 Compilation Cost}\label{compilation-cost}

Measure synthesis time as the number of handlers and precedence
constraints increases.

The expected complexity is:

\begin{lstlisting}
O(|H| + |≺|)
\end{lstlisting}

for the ordering phase.

The experiment should verify that synthesis remains negligible relative
to ordinary compilation and execution costs for realistic handler
counts.

\section{12. Relationship to Prior
Work}\label{relationship-to-prior-work}

\subsection{12.1 The Arrow Calculus}\label{the-arrow-calculus}

The arrow calculus of Lindley, Wadler, and Yallop provides the
foundational calculus for reasoning about computations with explicit
structure.

Our work does not introduce a new arrow calculus. Instead, it uses the
existing handler mechanism as the target of a declarative compilation
layer.

\subsection{12.2 Lindley's λ\_flow}\label{lindleys-ux3bb_flow}

Lindley's λ\_flow explored algebraic effects and handlers across idioms,
arrows, and monads and provided static effect distinctions.

This work is relevant because it demonstrates that effect structure can
be exposed statically.

Our contribution is different: we use a declarative specification to
\textbf{construct the handler topology itself}.

Thus the relationship can be summarized as:

\begin{lstlisting}
λ_flow
  │
  │ static effect structure
  ▼
Sanada
  │
  │ formal arrow handlers
  ▼
Declarative synthesis
  │
  │ handler precedence
  ▼
nested handle expressions
\end{lstlisting}

\subsection{12.3 Algebraic Effects and
Handlers}\label{algebraic-effects-and-handlers}

General algebraic-effects research establishes the separation between
effect declaration and interpretation.

Our system applies this principle specifically to handler composition:
the declarative specification describes how handlers should relate,
while the compiler determines the concrete nesting required by the
target calculus.

\section{13. Limitations}\label{limitations}

The system intentionally has a narrow scope.

First, it synthesizes a \textbf{static linear handler topology}. It does
not synthesize arbitrary runtime handler graphs.

Second, the current model does not attempt to infer handler precedence
automatically from arbitrary program behavior. Precedence must either be
declared or derived from a separately specified static policy.

Third, a partial order may admit multiple valid linear extensions.
Although deterministic tie-breaking can provide reproducible output,
different valid extensions may have different runtime behavior when
handlers interact. A complete system therefore needs a semantic account
of when two handlers may be safely reordered.

Fourth, the present work does not establish that every useful effect
system can be represented by the arrow model. Its claims are restricted
to handler topology that can be represented by the existing arrow
calculus.

Finally, the current implementation does not optimize handler
composition. Its purpose is to establish the correctness and
practicality of the declarative-to-nested transformation.

\section{14. Future Work}\label{future-work}

\subsection{14.1 Semantic Commutativity}\label{semantic-commutativity}

An important question is when two handlers can safely be reordered.

We provide a partial answer, formalized in
\passthrough{\lstinline!Commutativity.lean!}. Two handlers are
\textbf{clause-disjoint} when no operation name is handled by both:

\begin{lstlisting}
disjointClauses(H₁, H₂) := ∀ op. ¬(H₁ handles op ∧ H₂ handles op)
\end{lstlisting}

\textbf{Theorem 4 --- Two-Handler Commutativity.} For a static body
\passthrough{\lstinline!P!}, if \passthrough{\lstinline!H₁!} and
\passthrough{\lstinline!H₂!} are clause-disjoint, then:

\begin{lstlisting}
eval(synthesize(P, [H₁, H₂])) = eval(synthesize(P, [H₂, H₁]))
\end{lstlisting}

That is, the trace is identical regardless of which disjoint handler is
innermost.

The proof proceeds by showing that swapping two disjoint adjacent
handlers in the handler context does not change which handler first
catches any operation (\passthrough{\lstinline!firstHandler\_swap!}),
and then lifting this to all static command forms by induction
(\passthrough{\lstinline!evalWith\_swap!}).

A counterexample demonstrates that the disjointness condition is
necessary: when two handlers both handle the same operation, the
innermost handler catches it, so reordering changes the trace.

The general case --- arbitrary permutations of pairwise-disjoint handler
lists --- is left as future work. The two-handler case is fully proved.

\subsection{14.2 Inferred Precedence}\label{inferred-precedence}

Some precedence constraints could potentially be inferred from operation
signatures or handler dependencies.

For example, if a handler invokes an operation that another handler must
intercept, the compiler could derive a necessary nesting relationship.

\subsection{14.3 Optimization}\label{optimization}

A synthesis pass could eliminate redundant handlers, merge compatible
handlers, or specialize handler nesting when their operation sets are
statically known.

\subsection{14.4 Beyond Formalization}\label{beyond-formalization}

The Lean formalization (§7) establishes both structural and semantic
correctness of synthesis: well-formedness, precedence preservation,
staticity, completeness, existence of linear extensions, semantic
equivalence between synthesized and manually nested terms (§7.2), and
two-handler commutativity for disjoint handlers (§14.1). Further
formalization work could verify:

\begin{enumerate}
\def\labelenumi{\arabic{enumi}.}
\tightlist
\item
  The general commutativity conjecture: for a static body, if all
  handlers in a list are pairwise clause-disjoint, then any permutation
  of that list produces the same trace.
\item
  Semantic equivalence under a richer operational semantics that models
  continuation capture and resumption, beyond the trace semantics
  currently formalized.
\end{enumerate}

\subsection{14.5 Beyond Total Handler
Order}\label{beyond-total-handler-order}

Although the final arrow representation is linear, future work could
investigate richer declarative descriptions that compile into multiple
statically selected handler configurations. For example, a specification
might describe two alternative handler orderings guarded by a
compile-time condition (e.g.~debug vs production mode). The compiler
would then generate two separate nested \passthrough{\lstinline!handle!}
expressions and select between them using permitted term-level branching
(\passthrough{\lstinline!if M then N₁ else N₂!} where
\passthrough{\lstinline!N₁!} and \passthrough{\lstinline!N₂!} are terms,
not commands). This stays within the arrow calculus because the
selection is static or term-level --- it does not introduce runtime
branching on effectful commands. The key insight is that multiple static
topologies can coexist as long as the choice between them does not
depend on operation results.

\section{15. Conclusion}\label{conclusion}

We presented a declarative system for synthesizing handler topology in
an arrow calculus with algebraic effects and handlers.

The key idea is to separate \textbf{handler specification} from
\textbf{handler syntax}. A programmer declares a set of handlers and a
precedence relation. The compiler computes a linear extension of that
relation and lowers it to the nested \passthrough{\lstinline!handle!}
expressions required by the underlying calculus.

This transformation is fundamentally static. It does not introduce
runtime branching over effectful commands and therefore does not require
an extension of the arrow semantic model. The synthesis problem is
consequently reduced to a well-defined compile-time ordering problem.
Its correctness has been mechanically verified in Lean 4 (§7), covering
well-formedness, precedence preservation, staticity, completeness,
semantic equivalence, and two-handler commutativity for disjoint
handlers.

The approach also provides a minimal implementation path. In a language
with existing handler-expression AST nodes, the declarative system can
be implemented as an AST-to-AST transformation producing those nodes.
The existing runtime remains responsible for ordinary handler
evaluation.

The principal result is therefore not a new runtime effect mechanism,
but a new \textbf{declarative layer for constructing an existing
arrow-effect mechanism}. The boundary is precise:

\begin{lstlisting}
static handler topology
        + synthesizable

runtime effect-dependent command topology
        - not introduced by this system
\end{lstlisting}

This boundary turns a limitation of arrow effects into a useful design
constraint. Because the generated handler structure is linear and
statically known, its synthesis is simple, deterministic, and amenable
to formal verification.

The resulting system provides a small example of a broader principle:
when a programming language exposes a structurally constrained semantic
representation, a declarative layer can synthesize that representation
without changing the underlying runtime. In the present case, Sanada's
explicit handler topology becomes not merely a syntactic requirement,
but the target of a compile-time synthesis process.

\section{References}\label{references}

\begin{enumerate}
\def\labelenumi{\arabic{enumi}.}
\tightlist
\item
  Sanada, T. (2024). \emph{Algebraic effects and handlers for arrows.}
  Journal of Functional Programming, 34. doi:10.1017/S0956796824000066
\item
  Sanada, T. (2024). \emph{Algebraic effects and handlers for arrows.}
  PhD thesis, Kyoto University.
  https://www.kurims.kyoto-u.ac.jp/\textasciitilde tsanada/papers/phdthesis.pdf
\item
  Lindley, S. (2014). \emph{Algebraic effects and effect handlers for
  idioms and arrows.} In Proceedings of the 19th International
  Conference on Functional Programming (ICFP).
  doi:10.1145/2633628.2633636
\item
  Lindley, S., Wadler, P., and Yallop, J. (2010). \emph{The arrow
  calculus.} Journal of Functional Programming, 20(1), pp.~51--87.
\item
  Plotkin, G. and Pretnar, M. (2013). \emph{Handling algebraic effects.}
  Logical Methods in Computer Science, 9(4).
\item
  Hughes, J. (2000). \emph{Generalising monads to arrows.} Science of
  Computer Programming, 37(1--3), pp.~67--111.
\end{enumerate}

\end{document}
